\documentclass[11pt,oneside]{article}

\usepackage[T1]{fontenc}
\usepackage{lmodern}
\usepackage{microtype}
\usepackage[margin=1in,headheight=14pt]{geometry}
\usepackage{amsmath,amssymb,amsthm,mathtools}
\usepackage{booktabs,longtable,tabularx,array}
\usepackage{enumitem}
\usepackage{xcolor}
\usepackage{xurl}
\usepackage{hyperref}
\usepackage{listings}
\usepackage{fancyhdr}
\usepackage{titlesec}
\usepackage{needspace}

\definecolor{deepblue}{RGB}{24,75,125}
\definecolor{proofgreen}{RGB}{16,105,69}
\definecolor{openred}{RGB}{160,45,45}
\definecolor{computeviolet}{RGB}{100,55,145}
\definecolor{graytext}{RGB}{85,92,99}
\definecolor{codegray}{RGB}{248,248,248}
\definecolor{rulegray}{RGB}{105,110,118}

\hypersetup{
  colorlinks=true,
  linkcolor=deepblue,
  citecolor=proofgreen,
  urlcolor=deepblue,
  pdftitle={Cyclotomic Gap Quotients and Recurrent GCD Signatures for the Alaoglu--Erdos Conjecture},
  pdfauthor={OpenAI GPT-5.6 Pro},
  pdfsubject={A nonduplicative continuation report on the Alaoglu--Erdos integer-exponent conjecture}
}
\urlstyle{same}

\pagestyle{fancy}
\fancyhf{}
\lhead{Alaoglu--Erd\H{o}s: cyclotomic continuation}
\rhead{22 August 2026}
\cfoot{\thepage}
\renewcommand{\headrulewidth}{0.4pt}

\titleformat{\section}{\Large\bfseries}{\thesection}{0.75em}{}
\titleformat{\subsection}{\large\bfseries}{\thesubsection}{0.75em}{}
\titleformat{\subsubsection}{\normalsize\bfseries}{\thesubsubsection}{0.75em}{}

\setlength{\parindent}{0pt}
\setlength{\parskip}{0.55em}
\setlength{\emergencystretch}{3em}
\setlist[itemize]{leftmargin=2em,itemsep=0.25em,topsep=0.35em}
\setlist[enumerate]{leftmargin=2.25em,itemsep=0.25em,topsep=0.35em}

\newtheorem{theorem}{Theorem}[section]
\newtheorem{proposition}[theorem]{Proposition}
\newtheorem{lemma}[theorem]{Lemma}
\newtheorem{corollary}[theorem]{Corollary}
\newtheorem{conjecture}[theorem]{Conjecture}
\newtheorem{criterion}[theorem]{Proof criterion}
\newtheorem{question}[theorem]{Open question}
\theoremstyle{definition}
\newtheorem{definition}[theorem]{Definition}
\newtheorem{experiment}[theorem]{Exact experiment}
\newtheorem{warning}[theorem]{Caution}
\theoremstyle{remark}
\newtheorem{remark}[theorem]{Remark}

\newcommand{\N}{\mathbb N}
\newcommand{\Nzero}{\mathbb N_0}
\newcommand{\Z}{\mathbb Z}
\newcommand{\Q}{\mathbb Q}
\newcommand{\R}{\mathbb R}
\newcommand{\Gm}{\mathbb G_m}
\newcommand{\ord}{\operatorname{ord}}
\newcommand{\den}{\operatorname{den}}
\newcommand{\num}{\operatorname{num}}
\newcommand{\rad}{\operatorname{rad}}
\newcommand{\supp}{\operatorname{supp}}
\newcommand{\vp}{v_p}
\newcommand{\ee}{\mathrm e}
\newcommand{\PhiH}{\boldsymbol{\Phi}}
\newcommand{\AEconj}{Alaoglu--Erd\H{o}s}
\newcommand{\statusproved}{\textcolor{proofgreen}{\textbf{[proved here]}}}
\newcommand{\statusexternal}{\textcolor{deepblue}{\textbf{[external theorem]}}}
\newcommand{\statuscomputed}{\textcolor{computeviolet}{\textbf{[exact computation]}}}
\newcommand{\statusopen}{\textcolor{openred}{\textbf{[open bridge]}}}
\newcommand{\statusinherited}{\textcolor{graytext}{\textbf{[inherited from the unified report]}}}

\newenvironment{statusbox}[1]
{\begin{center}\begin{minipage}{0.94\textwidth}\hrule\vspace{0.45em}\textbf{Status:} #1\par\vspace{0.35em}}
{\vspace{0.35em}\hrule\end{minipage}\end{center}}

\newenvironment{resultbox}
{\begin{center}\begin{minipage}{0.94\textwidth}\hrule\vspace{0.5em}}
{\vspace{0.45em}\hrule\end{minipage}\end{center}}

\lstdefinestyle{pythonstyle}{
  language=Python,
  basicstyle=\ttfamily\scriptsize,
  backgroundcolor=\color{codegray},
  frame=single,
  rulecolor=\color{rulegray},
  numbers=left,
  numberstyle=\tiny\color{rulegray},
  numbersep=7pt,
  showstringspaces=false,
  breaklines=true,
  columns=fullflexible,
  keepspaces=true,
  tabsize=4
}

\begin{document}
\pagenumbering{gobble}

\begin{titlepage}
\centering
\vspace*{0.8cm}
{\Huge\bfseries Cyclotomic Gap Quotients,\par}
\vspace{0.15cm}
{\Huge\bfseries Recurrent GCD Signatures,\par}
\vspace{0.15cm}
{\Huge\bfseries and Finite-Field Synchronization\par}
\vspace{0.55cm}
{\LARGE A nonduplicative continuation of the unified report on the\par}
\vspace{0.1cm}
{\LARGE \AEconj{} integer-exponent conjecture\par}
\vspace{0.75cm}
{\large 22 August 2026\par}
\vfill
\begin{minipage}{0.91\textwidth}
\small
\textbf{Problem.}  Prove that, for real $x$,
\[
  2^x,3^x\in\Z \quad\Longrightarrow\quad x\in\Z.
\]

\textbf{Bottom line.}  I do not obtain a complete proof.  I do obtain a new, exact
cyclotomic dichotomy.  Fix any nonzero lattice direction $(a,b)\in\Z^2$ and reduce
$r=2^a3^b=A/B>1$ and $r^x=M^aN^b=C/D>1$.  If $x$ is a positive integer, then
\[
  \frac{C^n-D^n}{A^n-B^n}\in\Z\qquad(n\ge1).
\]
If $x$ is a hypothetical noninteger solution, then instead
\[
 \log\den\!\left(\frac{C^n-D^n}{A^n-B^n}\right)
   =n\log A+o(n).
\]
Thus essentially the entire exponential input gap survives in the reduced denominator,
in every rational direction at once.  This converts the conjecture into several precise
finite-place targets: any fixed direction with a positive exponential denominator saving,
any fixed finite prime support along infinitely many layers, any almost-everywhere support
inclusion, or any bounded synchronization cost would finish the proof.

A second new package imports a June 2026 theorem on the sequence
$\gcd(a^n-1,b^n-1)$.  It yields a recurrence-signature characterization of integrality and
shows exactly why constructing one nonperiodic common linear divisibility sequence would be
a proof certificate.  Finally, an exterior-power calculation shows that ordinary Hadamard
sign preconditioning is globally determinant-neutral; any benefit must come from selective
arithmetic structure, not orthogonal mixing itself.
\end{minipage}
\vfill
{\small Prepared from the attached unified report, but deliberately organized around
objects and theorems absent from that report.}
\end{titlepage}

\clearpage
\pagenumbering{roman}
\tableofcontents
\clearpage
\pagenumbering{arabic}

\section{Executive assessment}
\label{sec:executive}

\subsection{What is new in this continuation}

The attached unified report already develops an exceptionally broad collection of methods:
conditional free-monoid structure, primitive odd cores, finite verification, arbitrary-node
and mixed-semigroup determinants, $p$-adic transport ceilings, Smith profiles, Loewner and
Pad\'e systems, rational secants between distinct lattice directions, Sturmian coding,
canonical products, moment methods, and several exact sufficient conditions.  Repeating any
of those approaches would add volume rather than information.

This continuation therefore starts from three objects not developed there:
\begin{enumerate}[label=\textup{(\roman*)}]
\item the \emph{same-direction power quotient}
\[
 Q_{\mathbf a}(n)=\frac{C_{\mathbf a}^{\,n}-D_{\mathbf a}^{\,n}}
 {A_{\mathbf a}^{\,n}-B_{\mathbf a}^{\,n}},
\]
not a secant between two distinct exponent vectors;
\item the \emph{source--output gcd sequence}
\[
 g_{2,M}(n)=\gcd(2^n-1,M^n-1),
\]
and its constant-coefficient recurrence signature;
\item the \emph{capture cost} $\kappa_{\mathbf a}(n)$, the least multiplier of a
hypothetical solution needed to make its output reproduce the full finite-field torsion of
an input gap at level $n$.
\end{enumerate}

The resulting theorems use the Bugeaud--Corvaja--Zannier and Corvaja--Zannier gcd estimates in
a way not present in the source: rather than recording them as a general obstruction to
large gcds, we compute the exact exponential rate of the reduced denominator of every
candidate gap quotient.  The report also incorporates the June 2026 preprint of
Nguyen-Dang on recurrent gcd sequences, which postdates much of the standard literature and
is not treated in the attached manuscript.

\subsection{Claim ledger}

\begin{longtable}{@{}p{0.29\textwidth}p{0.19\textwidth}p{0.45\textwidth}@{}}
\toprule
Statement & Status & Dependency \\
\midrule
\endfirsthead
\toprule
Statement & Status & Dependency \\
\midrule
\endhead
Directional gap-quotient dichotomy & \statusproved & Corvaja--Zannier's generalized
$S$-unit gcd theorem, plus elementary height and prime-valuation arguments. \\
Recurrence signature of $g_{2,M}$ & \statusexternal & Nguyen-Dang, arXiv:2606.07959
(2026), itself using the classical BCZ estimate and linear-divisibility-sequence theory. \\
Primitive-core collapse of the diagonal gcd sequences & \statusproved & Elementary
congruences modulo $2^n-1$ and $3^n-1$. \\
Fixed-support, denominator-saving, and multiplier criteria & \statusproved & Immediate
consequences of the directional dichotomy and LTE. \\
Support-inclusion criterion & \statusproved & Corrales-Rodrig\'a\~nez--Schoof's support
problem theorem. \\
Exact synchronization-cost formula and divergence & \statusproved & Multiplicative
orders, Euler's theorem, and the homogeneous gcd estimate. \\
Numerical tables & \statuscomputed & Exact integer factorization and multiplicative orders;
reproducible code is included. \\
Hadamard preconditioning neutrality & \statusproved & Orthogonal invariance of singular
values and exterior powers. \\
Derivation of bounded synchronization from $2^x,3^x\in\Z$ & \statusopen & This is the
newly isolated archimedean-to-Frobenius bridge. \\
Complete proof of the conjecture & \statusopen & No contradiction with a hypothetical
counterexample has been obtained. \\
\bottomrule
\end{longtable}

\subsection{The strongest theorem in one line}

\begin{resultbox}
For a positive solution $x$ and every nonzero $\mathbf a=(a,b)\in\Z^2$, orient
$r_{\mathbf a}=2^a3^b=A/B>1$ and reduce
$r_{\mathbf a}^{x}=M^aN^b=C/D>1$.  Then
\[
 x\in\Z
 \quad\Longrightarrow\quad
 \den Q_{\mathbf a}(n)=1\quad(n\ge1),
\]
whereas
\[
 x\notin\Z
 \quad\Longrightarrow\quad
 \frac{\log\den Q_{\mathbf a}(n)}{n}\longrightarrow\log A.
\]
No intermediate asymptotic regime is possible.
\end{resultbox}

This is not merely another equivalent formulation.  It gives a quantitative rigidity law:
a counterexample must be maximally decorrelated from the input at finite places, despite
being perfectly correlated at the real place by the identity
$M^aN^b=(2^a3^b)^x$.

\subsection{External-claims audit}

The public mathematical record was checked because the prompt reports several very recent
AI-assisted breakthroughs.  The Jacobian claim has a publicly displayed explicit
counterexample and substantial secondary discussion.  The Hadamard claim is less settled:
secondary posts assert that twelve formerly open orders below $2000$ were filled, but current
public Sage/open-search records still list order $668$ as unresolved and I did not locate the
claimed exact matrices.  The rank-$30$ elliptic-curve claim is visible in social mirrors and
aggregators, but I did not locate a primary paper or independent certificate with settled
attribution.  Most importantly for the present task, I found no public preprint, Lean
repository, or checkable argument for a new \AEconj{} breakthrough as of 22 August 2026.
Accordingly, no technical step from the rival team's private claim is assumed here.

\begin{statusbox}{The news audit is evidentiary context, not a mathematical premise.  The
Hadamard and rank-$30$ claims should be treated as provisional until exact public
certificates are available.}
\end{statusbox}

\section{Minimal inherited setup}
\label{sec:setup}

Let $x\in\R$ satisfy
\[
 M:=2^x\in\Z,\qquad N:=3^x\in\Z.
\]
If $x<0$, then $0<2^x<1$, so no solution is possible.  The case $x=0$ is integral and
trivial.  Hence every counterexample would have $x>0$ and $M,N\ge2$.

We use only the following inherited facts from the unified report:
\begin{enumerate}[label=\textup{(\roman*)}]
\item a rational solution is an integer, by unique factorization;
\item every nonintegral solution is transcendental;
\item positive integer multiples of a solution are again solutions;
\item conditional on failure, one may choose a least nonintegral solution $\beta$ with a
primitive output normalization
\[
  M=2^\beta=2^a w^d,
\]
where $w>1$ is odd and power-primitive; symmetrically
$N=3^b v^e$, and $a=0$ or $b=0$;
\item the full conditional solution set is the free monoid
$\{j+k\beta:j,k\in\Nzero\}$.
\end{enumerate}

Only items (i) and (iii) are needed for most of this report.  The primitive normalization is
used to show that the new gcd and synchronization invariants ignore the base-adic depth and
see only the primitive core.

\begin{lemma}[Rational-exponent lemma]
\label{lem:rational}
If $q\in\Q$ and $2^q\in\Z$, then $q\in\Z$.
\end{lemma}

\begin{proof}
Write $q=r/s$ in lowest terms with $s>0$.  If $2^{r/s}=m\in\Z$, then
$m^s=2^r$.  Unique factorization gives $m=2^t$ and $ts=r$.  Since
$\gcd(r,s)=1$, we have $s=1$.
\end{proof}

\begin{lemma}[Output independence]
\label{lem:MNindependent}
For every positive solution $x$, the integers $M$ and $N$ are multiplicatively independent.
\end{lemma}

\begin{proof}
If $M^r=N^s$ for positive integers $r,s$, then
$2^{rx}=3^{sx}$.  Since $x>0$, this gives $2^r=3^s$, impossible.
\end{proof}

\section{Recurrent gcd signatures}
\label{sec:recurrence}

\subsection{The 2026 recurrence theorem}

For integers $a,b\ge2$, put
\[
 g_{a,b}(n):=\gcd(a^n-1,b^n-1)\qquad(n\ge1).
\]
A sequence is \emph{$C$-finite} if it satisfies a constant-coefficient linear recurrence.
It is a \emph{linear divisibility sequence} (LDS) if it is both $C$-finite and satisfies
$W_m\mid W_n$ whenever $m\mid n$.

\begin{theorem}[Nguyen-Dang, 2026]
\label{thm:ND}
Let $a,b\ge2$.  The following are equivalent:
\begin{enumerate}[label=\textup{(\alph*)}]
\item $a$ and $b$ are multiplicatively dependent;
\item $(g_{a,b}(n))_{n\ge1}$ is $C$-finite;
\item some tail of $(g_{a,b}(n))$ is $C$-finite.
\end{enumerate}
If $a,b$ are multiplicatively independent, every integer LDS $W_n$ satisfying
\[
 W_n\mid a^n-1,\qquad W_n\mid b^n-1\qquad(n\ge1)
\]
is periodic.
\end{theorem}

\begin{statusbox}{\statusexternal{} --- Theorem 1.1 and Theorem 1.2 of
Nguyen-Dang~\cite{NguyenDang2026}.  The preprint is recent and should be cited as a preprint,
not as a peer-reviewed theorem.}
\end{statusbox}

\subsection{A recurrence certificate for integrality}

\begin{theorem}[One-base recurrence signature]
\label{thm:onebase-signature}
Let $x>0$ and suppose $M=2^x\in\Z$.  The following are equivalent:
\begin{enumerate}[label=\textup{(\roman*)}]
\item $x\in\Z$;
\item $2$ and $M$ are multiplicatively dependent;
\item $g_{2,M}(n)$ is $C$-finite;
\item some tail of $g_{2,M}(n)$ is $C$-finite;
\item there exists a nonperiodic integer LDS $W_n$ such that
$W_n\mid2^n-1$ and $W_n\mid M^n-1$ for all $n$.
\end{enumerate}
The analogous equivalence holds with $(3,N)$ in place of $(2,M)$.
\end{theorem}

\begin{proof}
If $2$ and $M$ are dependent, then $M^s=2^r$ for some $r,s>0$.
Unique factorization forces $M=2^k$ for an integer $k$, hence $x=k$.
The converse is immediate.  The equivalence with (iii)--(iv) is
Theorem~\ref{thm:ND}.  If $x=k\in\Z_{>0}$, then
$W_n=2^n-1$ is a nonperiodic LDS and divides
$M^n-1=2^{kn}-1$.  If $x\notin\Z$, Theorem~\ref{thm:ND} says that every such common LDS is
periodic.
\end{proof}

\begin{criterion}[LDS proof certificate]
\label{crit:LDS}
To prove the \AEconj{} conjecture it would suffice, from the simultaneous hypotheses
$2^x,3^x\in\Z$, to construct a single nonperiodic LDS dividing both
$2^n-1$ and $(2^x)^n-1$ for every $n$.
\end{criterion}

This criterion is much more rigid than merely constructing a divisibility sequence: the gcd
sequence itself is always a divisibility sequence.  The missing property is finite linear
recurrence, or a nonperiodic recurrent factor.

\subsection{The recurrence graph of the four integers}

For $x>0$ define a graph $\mathcal R_x$ on the vertices
\[
 \{2,3,M,N\}
\]
by joining $a$ to $b$ exactly when $g_{a,b}(n)$ is $C$-finite.  By
Theorem~\ref{thm:ND}, this is exactly the multiplicative-dependence graph.

\begin{theorem}[Recurrence-graph phase transition]
\label{thm:graph}
Let $x>0$ satisfy $M=2^x,N=3^x\in\Z$.
\begin{enumerate}[label=\textup{(\roman*)}]
\item If $x\in\Z$, then $\mathcal R_x$ has exactly the two edges
$\{2,M\}$ and $\{3,N\}$.
\item If $x\notin\Z$, then neither of those two edges occurs.  The only possible edges are
$\{2,N\}$ and $\{3,M\}$, and at most one of them can occur.
\end{enumerate}
\end{theorem}

\begin{proof}
The source pair $(2,3)$ and output pair $(M,N)$ are independent, the latter by
Lemma~\ref{lem:MNindependent}.  The pairs $(2,M)$ and $(3,N)$ are dependent exactly when
$x$ is integral, by Theorem~\ref{thm:onebase-signature}.

It remains to exclude simultaneous cross-dependence in the nonintegral case.  If both cross
edges occurred, unique factorization would give $N=2^u$ and $M=3^v$ for some positive
integers $u,v$.  With $\vartheta=\log_2 3$, the equations
$3^x=2^u$ and $2^x=3^v$ give $x\vartheta=u$ and $x=v\vartheta$, hence
$\vartheta^2=u/v\in\Q$.  But $\vartheta$ is transcendental: rationality is impossible by
unique factorization, while algebraic irrationality would contradict Gelfond--Schneider
because $2^\vartheta=3$.  Thus both cross edges cannot occur.
\end{proof}

\begin{remark}
The graph theorem is a clean diagnostic, but not a proof: a generic hypothetical
counterexample can simply have no recurrent edges.  Its value is to state exactly what a
successful recurrence construction must change.
\end{remark}

\subsection{Primitive-core collapse}

The conditional normal form in the unified report becomes especially transparent modulo
$2^n-1$.

\begin{lemma}[Dyadic depth is invisible]
\label{lem:core-collapse}
If $M=2^aR$ with $a\ge0$, then for every $n\ge1$,
\[
 \gcd(2^n-1,M^n-1)=\gcd(2^n-1,R^n-1).
\]
In particular, if a least counterexample has $M=2^aw^d$, then
\[
 g_{2,M}(n)=\gcd(2^n-1,w^{dn}-1).
\]
Similarly, if $N=3^bS$, then
\[
 g_{3,N}(n)=\gcd(3^n-1,S^n-1).
\]
\end{lemma}

\begin{proof}
Modulo $2^n-1$ one has $2^n\equiv1$, so
$M^n=2^{an}R^n\equiv R^n$.  The gcd identity follows.  The base-$3$ statement is identical.
\end{proof}

Thus the new diagonal gcd sequence forgets the base-adic depth $a$ or $b$ and sees only the
primitive base-free core.  This explains, for example, why the exact sequences for
$M=3,6,12$ coincide.

\section{Homogeneous gcds and directional gap quotients}
\label{sec:directional}

\subsection{The homogeneous $S$-unit gcd estimate}

We need a rational version of the classical BCZ estimate.

\begin{theorem}[Homogeneous gcd estimate]
\label{thm:homogeneous-gcd}
Let $A>B\ge1$ and $C>D\ge1$ be coprime pairs.  If $A/B$ and $C/D$ are
multiplicatively independent, then
\[
 \log\gcd(A^n-B^n,C^n-D^n)=o(n).
\]
Equivalently, for every $\varepsilon>0$ the gcd is at most $\exp(\varepsilon n)$ for all
sufficiently large $n$.
\end{theorem}

\begin{proof}[Derivation from Corvaja--Zannier]
Apply Corollary 1 of Corvaja--Zannier~\cite{CorvajaZannier2005} to the cyclic subgroup
\[
 \left\{\bigl((A/B)^n,(C/D)^n\bigr):n\in\Z\right\}\subset\Gm^2(\Q).
\]
Multiplicative independence keeps this subgroup outside every exceptional one-dimensional
subtorus relevant to $(u-1,v-1)$.  Their generalized gcd estimate gives sublinear logarithmic
common finite-place mass.  Because $(A,B)=(C,D)=1$, a prime dividing $A$ or $B$ cannot divide
$A^n-B^n$, and similarly for $C,D$.  Hence every prime in the ordinary integer gcd is a
nonexceptional finite place, and the generalized estimate is exactly the displayed
$o(n)$ bound.
\end{proof}

\begin{statusbox}{\statusexternal{} for the Subspace-Theorem input;
\statusproved{} for the homogeneous specialization and all applications below.  The original
integer case $B=D=1$ is the theorem of Bugeaud--Corvaja--Zannier~\cite{BCZ2003}.}
\end{statusbox}

\subsection{A rational direction attached to a solution}

Let $\mathbf a=(a,b)\in\Z^2\setminus\{(0,0)\}$.  Since $2$ and $3$ are
multiplicatively independent, $2^a3^b\ne1$.  Replace $\mathbf a$ by $-\mathbf a$ if needed
and write in lowest terms
\[
 r_{\mathbf a}:=2^a3^b=\frac{A_{\mathbf a}}{B_{\mathbf a}}>1,
 \qquad (A_{\mathbf a},B_{\mathbf a})=1.
\]
For a positive solution $x$ define
\[
 s_{\mathbf a}:=M^aN^b=r_{\mathbf a}^{\,x}
 =\frac{C_{\mathbf a}}{D_{\mathbf a}}>1,
 \qquad(C_{\mathbf a},D_{\mathbf a})=1.
\]
The identity $s_{\mathbf a}=r_{\mathbf a}^x$ follows simply from
$M=2^x,N=3^x$ and remains valid for negative $a$ or $b$.

Define the homogeneous gaps and their quotient by
\[
 U_{\mathbf a}(n)=A_{\mathbf a}^{\,n}-B_{\mathbf a}^{\,n},
 \qquad
 V_{\mathbf a}(n)=C_{\mathbf a}^{\,n}-D_{\mathbf a}^{\,n},
 \qquad
 Q_{\mathbf a}(n)=\frac{V_{\mathbf a}(n)}{U_{\mathbf a}(n)}.
\]
Also put
\[
 G_{\mathbf a}(n)=\gcd(U_{\mathbf a}(n),V_{\mathbf a}(n)),
 \quad
 \Delta_{\mathbf a}(n)=\den Q_{\mathbf a}(n)
 =\frac{U_{\mathbf a}(n)}{G_{\mathbf a}(n)}.
\]

\begin{lemma}[Direction dependence detects rationality]
\label{lem:direction-dependence}
For any nonzero direction $\mathbf a$, the rationals $r_{\mathbf a}$ and
$s_{\mathbf a}$ are multiplicatively dependent if and only if $x\in\Q$.  Under the
hypothesis $M=2^x\in\Z$, this is equivalent to $x\in\Z$.
\end{lemma}

\begin{proof}
If $r^u=s^v$ with positive integers $u,v$, then
$r^u=(r^x)^v$ and $r>1$, so $x=u/v$.  Conversely, if $x=u/v$, then $s^v=r^u$.
The last assertion is Lemma~\ref{lem:rational}.
\end{proof}

\subsection{The directional denominator dichotomy}

\begin{theorem}[Directional gap-quotient dichotomy]
\label{thm:directional-dichotomy}
Let $x>0$ satisfy $M=2^x,N=3^x\in\Z$, and fix any nonzero direction
$\mathbf a\in\Z^2$ oriented as above.
\begin{enumerate}[label=\textup{(\roman*)}]
\item If $x=k\in\Z_{>0}$, then
\[
 Q_{\mathbf a}(n)
 =\frac{A_{\mathbf a}^{kn}-B_{\mathbf a}^{kn}}
 {A_{\mathbf a}^{n}-B_{\mathbf a}^{n}}
 \in\Z
 \qquad(n\ge1).
\]
In particular, $\Delta_{\mathbf a}(n)=1$ for every $n$.
\item If $x\notin\Z$, then
\[
 \log G_{\mathbf a}(n)=o(n),
\]
and
\begin{align*}
 \log \Delta_{\mathbf a}(n)&=n\log A_{\mathbf a}+o(n),\\
 \log \num Q_{\mathbf a}(n)&=n\log C_{\mathbf a}+o(n),\\
 h(Q_{\mathbf a}(n))&=
 n\max\{\log A_{\mathbf a},\log C_{\mathbf a}\}+o(n).
\end{align*}
\end{enumerate}
\end{theorem}

\begin{proof}
If $x=k\in\Z_{>0}$, reduction gives
$C_{\mathbf a}=A_{\mathbf a}^k$ and
$D_{\mathbf a}=B_{\mathbf a}^k$.  The standard factorization
$X^k-Y^k=(X-Y)(X^{k-1}+\cdots+Y^{k-1})$, with
$X=A_{\mathbf a}^{n}$ and $Y=B_{\mathbf a}^{n}$, proves (i).

Suppose $x\notin\Z$.  By Lemma~\ref{lem:direction-dependence}, the ratios
$A_{\mathbf a}/B_{\mathbf a}$ and $C_{\mathbf a}/D_{\mathbf a}$ are independent.
Theorem~\ref{thm:homogeneous-gcd} gives $\log G_{\mathbf a}(n)=o(n)$.  Since
\[
 \log U_{\mathbf a}(n)
 =n\log A_{\mathbf a}
 +\log\left(1-(B_{\mathbf a}/A_{\mathbf a})^n\right)
 =n\log A_{\mathbf a}+o(1),
\]
we obtain
\[
 \log\Delta_{\mathbf a}(n)
 =\log U_{\mathbf a}(n)-\log G_{\mathbf a}(n)
 =n\log A_{\mathbf a}+o(n).
\]
The numerator and height formulas are identical.
\end{proof}

\begin{corollary}[One direction is enough]
\label{cor:one-direction}
For a positive solution $x$ and any fixed nonzero direction $\mathbf a$, the following are
equivalent:
\begin{enumerate}[label=\textup{(\alph*)}]
\item $x\in\Z$;
\item $Q_{\mathbf a}(n)\in\Z$ for every $n\ge1$;
\item $Q_{\mathbf a}(n)\in\Z$ for infinitely many unbounded $n$;
\item $\log\Delta_{\mathbf a}(n)=o(n)$ along an unbounded sequence of $n$;
\item for some $\delta>0$,
$\Delta_{\mathbf a}(n)\le\exp((\log A_{\mathbf a}-\delta)n)$ for infinitely many $n$.
\end{enumerate}
\end{corollary}

\begin{proof}
The integral case follows from Theorem~\ref{thm:directional-dichotomy}(i).  Every one of
(b)--(e) contradicts the limiting denominator rate in part (ii) unless $x$ is integral.
\end{proof}

\begin{remark}[Why this does not yet prove the conjecture]
The original hypotheses give one exact integral value in some directions, but not infinitely
many.  For example,
\[
 \frac{M-1}{2-1}=M-1\in\Z.
\]
The theorem says that a noninteger solution can accommodate any finite collection of such
exceptions; the obstruction is asymptotic.
\end{remark}

\subsection{The four adjacent smooth directions}

The prior report isolates the consecutive $3$-smooth pairs
$(1,2),(2,3),(3,4),(8,9)$.  In the new language they correspond to the directions
\[
 2,\qquad \frac32,\qquad\frac43,\qquad\frac98.
\]
Their output ratios are respectively
\[
 M,\qquad\frac NM,\qquad\frac{M^2}{N},\qquad\frac{N^2}{M^3}.
\]
After reducing each output ratio to $C/D$, the denominator of
$(C-D)/(A-B)$ is $1$ because every input gap $A-B$ is $1$.

\begin{corollary}[One-step integrality followed by denominator explosion]
\label{cor:adjacent}
If $x$ is a hypothetical noninteger solution, then for each of the four directions above
one has $Q(1)\in\Z$, but
\[
 \log\den Q(n)=
 \begin{cases}
 n\log2+o(n),&r=2,\\
 n\log3+o(n),&r=3/2,\\
 n\log4+o(n),&r=4/3,\\
 n\log9+o(n),&r=9/8.
 \end{cases}
\]
\end{corollary}

This is a useful calibration: exact integrality at the first cyclotomic layer is automatic,
while nearly maximal denominator growth is forced at high layers.  Any proof based on the
adjacent smooth pairs must therefore propagate $n=1$ information to an unbounded set of
power layers; local convexity or first-order gap information alone cannot do that.


\section{Denominator-saving and cyclotomic criteria}
\label{sec:criteria}

The asymptotic formula in Theorem~\ref{thm:directional-dichotomy} becomes useful only after
it is translated into conditions that a prospective argument might actually produce.  This
section records several such translations.  They are elementary once the homogeneous gcd
estimate is available, but they sharply delimit what remains to be proved.

\subsection{A uniform valuation bound at fixed primes}

\begin{lemma}[Fixed-prime valuation bound]
\label{lem:fixed-prime-valuation}
Let $A>B\ge1$ be coprime integers and let $p$ be a fixed prime.  Then
\[
 v_p(A^n-B^n)=O_{A,B,p}(\log(n+1)).
\]
If $p\mid AB$, the valuation is identically zero.
\end{lemma}

\begin{proof}
If $p\mid A$, then $p\nmid B$ and $A^n-B^n\not\equiv0\pmod p$; the same argument applies
when $p\mid B$.  Assume now $p\nmid AB$.  If no power $(A/B)^n$ is $1$ modulo $p$, the
valuation is again zero.  Otherwise let $f=\ord_p(A/B)$.  A nonzero valuation requires
$f\mid n$.  For odd $p$, LTE gives, after fixing the first admissible exponent $f$,
\[
 v_p(A^n-B^n)=v_p(A^f-B^f)+v_p(n/f).
\]
For $p=2$, the standard two-adic LTE formula has the same form up to an additive constant
depending only on $A$ and $B$.  Since $v_p(n)\le \log n/\log p$, the assertion follows.
\end{proof}

The lemma says that a fixed finite set of primes can carry only polynomial, never
exponential, mass from a binary power gap.

\subsection{Equivalent one-direction proof targets}

Fix a direction $\mathbf a$ and abbreviate
\[
 U_n=A^n-B^n,\qquad V_n=C^n-D^n,\qquad
 G_n=\gcd(U_n,V_n),\qquad \Delta_n=U_n/G_n.
\]
Here $A/B=2^a3^b>1$ and $C/D=(A/B)^x$ are in lowest terms.

\begin{theorem}[A menu of finite-place proof criteria]
\label{thm:criteria-menu}
Suppose $x>0$ and $2^x,3^x\in\Z$.  If, for one fixed nonzero direction, any one of the
following assertions holds along an unbounded sequence of indices $n$, then $x\in\Z$:
\begin{enumerate}[label=\textup{(\alph*)}]
\item there is a fixed finite set of primes $S$ with $\supp(\Delta_n)\subseteq S$;
\item there are positive integers $K_n$ with
      $U_n\mid K_nV_n$ and $\log K_n=o(n)$;
\item there is a $\delta>0$ with $G_n\ge \exp(\delta n)$;
\item there is a $\delta>0$ with
      $\Delta_n\le \exp((\log A-\delta)n)$;
\item the radical satisfies
      $\log\rad(\Delta_n)=o(n/\log n)$ and the exponents
      $v_p(\Delta_n)$ are uniformly $O(\log n)$ for primes in its support.
\end{enumerate}
\end{theorem}

\begin{proof}
Assume for contradiction that $x\notin\Z$.  Then
$\log\Delta_n=n\log A+o(n)$ and $\log G_n=o(n)$.

For (a), Lemma~\ref{lem:fixed-prime-valuation} gives
\[
 \log\Delta_n
 \le \sum_{p\in S}v_p(U_n)\log p
 =O_S(\log n),
\]
a contradiction.  For (b), the divisibility $U_n\mid K_nV_n$ implies
$\Delta_n\mid K_n$: after cancelling $G_n$, the two integers
$U_n/G_n$ and $V_n/G_n$ are coprime.  Hence
$\log\Delta_n\le\log K_n=o(n)$, again impossible.  Conditions (c) and (d) directly
contradict the two asymptotic formulas.  Finally, (e) gives
\[
 \log\Delta_n
 =\sum_{p\mid\Delta_n}v_p(\Delta_n)\log p
 \le O(\log n)\sum_{p\mid\Delta_n}\log p
 =o(n),
\]
again contradicting denominator saturation.
\end{proof}

\begin{warning}[Radicals alone are delicate]
The bare estimate $\log\rad(\Delta_n)=o(n)$ is not by itself enough: a small collection of
primes whose identities vary with $n$ could carry large powers.  A successful radical
argument must control multiplicities as well.  The fixed-prime lemma supplies such control
only after the prime set has been stabilized.
\end{warning}

\begin{corollary}[Subexponential correction is fatal]
\label{cor:subexp-correction}
For a hypothetical counterexample and every fixed direction,
\[
 A^n-B^n\nmid K_n(C^n-D^n)
\]
for all sufficiently large $n$ whenever $K_n=\exp(o(n))$ is integral.
\end{corollary}

This is a useful target for approximation arguments.  One need not force exact divisibility
$U_n\mid V_n$; a multiplier of polynomial size, or even $\exp(o(n))$ size, would already
contradict a counterexample.

\subsection{Homogeneous cyclotomic layers}

Write $\Phi_m(X,Y)$ for the homogeneous cyclotomic polynomial, characterized by
\[
 X^n-Y^n=\prod_{m\mid n}\Phi_m(X,Y).
\]
Thus $U_n$ is assembled from layers indexed by divisors of $n$.  The denominator
$\Delta_n=U_n/G_n$ measures exactly the prime-power mass of those layers that is not
reproduced by the output gap $V_n$.

\begin{proposition}[Collision-mass identity]
\label{prop:collision-mass}
For every $n$,
\begin{align*}
 \log G_n
 &=\sum_p\min\{v_p(U_n),v_p(V_n)\}\log p,\\
 \log\Delta_n
 &=\sum_p\bigl(v_p(U_n)-\min\{v_p(U_n),v_p(V_n)\}\bigr)\log p.
\end{align*}
If $x\notin\Z$, then the first sum is $o(n)$ and the second is
$n\log A+o(n)$.
\end{proposition}

\begin{proof}
The identities are the prime factorizations of $G_n$ and $U_n/G_n$; the asymptotics are
Theorem~\ref{thm:directional-dichotomy}.
\end{proof}

The proposition should be read as a conservation law.  The input has total logarithmic mass
$n\log A+o(n)$.  In a counterexample, only $o(n)$ of that mass can collide with the output;
almost all of it survives in the denominator.

\begin{theorem}[Zsigmondy]
\label{thm:zsigmondy}
Let $A>B\ge1$ be coprime.  Except for the classical cases
\begin{enumerate}[label=\textup{(\roman*)}]
\item $n=2$ and $A+B$ is a power of $2$;
\item $n=6$ and $(A,B)=(2,1)$,
\end{enumerate}
there is, for every $n>1$, a prime dividing $A^n-B^n$ that divides no earlier
$A^m-B^m$ with $1\le m<n$.
\end{theorem}

For the adjacent direction $A/B=3/2$, neither exception occurs.  Consequently every layer
$n\ge2$ introduces a genuinely new input prime.  This observation alone is not enough:
the output is allowed to acquire the same new prime at the same layer.  The point of the
gcd theorem is that such synchronized acquisitions can carry only $o(n)$ total logarithmic
mass when the two ratios are independent.

\begin{corollary}[Prime-index cyclotomic saturation]
\label{cor:prime-layer}
Let $\ell$ tend to infinity through primes.  If $x\notin\Z$, then
\[
 \log\den\!\left(\frac{C^\ell-D^\ell}{A^\ell-B^\ell}\right)
   =\ell\log A+o(\ell).
\]
Since
\[
 A^\ell-B^\ell=(A-B)\Phi_\ell(A,B),
\]
the fixed factor $A-B$ contributes only $O(1)$; asymptotically all of the denominator mass
comes from the single new prime-index cyclotomic layer $\Phi_\ell(A,B)$.
\end{corollary}

\begin{remark}[Why primitive divisors do not close the proof]
Zsigmondy's theorem guarantees at least one fresh prime, not a large fresh factor.  The fresh
prime could be small compared with $A^n$.  To finish by primitive divisors one needs a
quantitative statement forcing a positive proportion of $\log\Phi_n(A,B)$ into primes or
prime powers that the output cannot synchronize.  Corollary~\ref{cor:prime-layer} proves the
aggregate failure of synchronization under the counterexample hypothesis, but it does not
connect that failure back to the original integrality at level $n=1$.
\end{remark}

\subsection{A useful contrapositive formulation}

\begin{criterion}[Cyclotomic propagation criterion]
\label{crit:cyclotomic-propagation}
It would prove the \AEconj{} conjecture to derive, from $M=2^x,N=3^x\in\Z$, one direction
$A/B$ and an unbounded sequence $n_j$ for which a positive fraction of the input gap is
reproduced by the output:
\[
 \log\gcd(A^{n_j}-B^{n_j},C^{n_j}-D^{n_j})\ge\delta n_j
\]
for some fixed $\delta>0$.
\end{criterion}

The desired fraction may be arbitrarily small.  This is a substantially softer target than
full divisibility, but it is still a genuinely finite-field statement: real approximation
of $C/D=(A/B)^x$ does not by itself imply any common prime factors.

\section{The support problem and unavoidable defect primes}
\label{sec:support}

The denominator dichotomy measures common prime-power mass.  A logically independent route
asks only which primes occur, ignoring valuations.  A theorem of
Corrales-Rodrig\'a\~nez and Schoof makes an almost-everywhere support inclusion completely
rigid.

\subsection{Support inclusion forces a power relation}

For a reduced rational number $r=A/B>1$ and a prime $p\nmid AB$, let
$\ord_p(r)$ denote the multiplicative order of $AB^{-1}$ in $\mathbb F_p^\times$.  Then
\[
 p\mid A^n-B^n \quad\Longleftrightarrow\quad \ord_p(r)\mid n.
\]
Thus inclusion between supports of all power gaps is equivalent to an order-divisibility
condition at almost every prime.

\begin{theorem}[Corrales--Rodrig\'a\~nez--Schoof]
\label{thm:support-problem}
Let $r,s\in\Q^\times$ be nontorsion.  Suppose that, for every positive integer $n$ and all
but finitely many primes $p$ at which both numbers are units,
\[
 r^n\equiv1\pmod p \quad\Longrightarrow\quad s^n\equiv1\pmod p.
\]
Then $s=r^k$ for some integer $k$.
\end{theorem}

\begin{statusbox}{\statusexternal{} --- this is the multiplicative-group case of the
support problem theorem~\cite{CorralesSchoof1997}.}
\end{statusbox}

\begin{theorem}[Support-inclusion criterion for \AEconj]
\label{thm:support-AE}
Fix one nonzero direction and write $r=A/B>1$, $s=C/D=r^x>1$.  Either one of the following
conditions implies $x\in\Z$:
\begin{enumerate}[label=\textup{(\alph*)}]
\item for every $n\ge1$ and all but finitely many primes $p$,
\[
 p\mid A^n-B^n\quad\Longrightarrow\quad p\mid C^n-D^n;
\]
\item for every $n\ge1$ and all but finitely many primes $p$,
\[
 p\mid C^n-D^n\quad\Longrightarrow\quad p\mid A^n-B^n.
\]
\end{enumerate}
\end{theorem}

\begin{proof}
Condition (a) and Theorem~\ref{thm:support-problem} give $s=r^k$ for an integer $k$.
Because $s=r^x$ and $r>1$, one has $x=k$.  Under (b), the theorem gives $r=s^k$, hence
$x=1/k\in\Q$.  Lemma~\ref{lem:rational} then yields $x\in\Z$ (in fact $x=1$ in this
orientation).
\end{proof}

In particular, it is enough to prove the inclusion for the diagonal pair $(2,M)$ or
$(3,N)$.  No information from both bases is needed once almost-all-prime support containment
has been obtained.

\subsection{Defect sets forced by a counterexample}

Define the directional defect sets
\begin{align*}
 \mathcal D_{r\to s}
 &=\left\{p:\text{ for some }n\ge1,
          \ p\mid A^n-B^n,\ p\nmid C^n-D^n\right\},\\
 \mathcal D_{s\to r}
 &=\left\{p:\text{ for some }n\ge1,
          \ p\mid C^n-D^n,\ p\nmid A^n-B^n\right\},
\end{align*}
with the finitely many primes dividing $ABCD$ omitted.

\begin{corollary}[Two-sided infinite support defects]
\label{cor:defect-infinite}
If $x$ is a hypothetical noninteger solution, then for every nonzero direction both
$\mathcal D_{r\to s}$ and $\mathcal D_{s\to r}$ are infinite.
\end{corollary}

\begin{proof}
If $\mathcal D_{r\to s}$ were finite, excluding it would give the support inclusion in
Theorem~\ref{thm:support-AE}(a), hence $x\in\Z$.  The reverse direction is identical.
\end{proof}

This is stronger than saying the two gap sequences are not equal.  Every rational direction
must generate infinitely many distinct primes witnessing failure of support inclusion in
each orientation.

\begin{proposition}[Order-theoretic form of the defect]
\label{prop:order-defect}
For $p\nmid ABCD$, the prime $p$ belongs to $\mathcal D_{r\to s}$ exactly when
\[
 \ord_p(s)\nmid \ord_p(r).
\]
Similarly, $p\in\mathcal D_{s\to r}$ exactly when
$\ord_p(r)\nmid\ord_p(s)$.
\end{proposition}

\begin{proof}
Take $n=\ord_p(r)$.  Then $r^n\equiv1\pmod p$.  The same exponent kills $s$ precisely
when $\ord_p(s)\mid\ord_p(r)$.  This proves the first equivalence, and the second follows by
symmetry.
\end{proof}

Consequently a counterexample must make the two order functions incomparable at infinitely
many primes in both directions.  The real identity $s=r^x$ imposes perfect proportionality
on logarithms, but the reductions of $r$ and $s$ modulo primes must repeatedly destroy every
integer divisibility relation between their orders.

\subsection{A support certificate with a finite exceptional set}

\begin{criterion}[Finite-exception support certificate]
\label{crit:support-certificate}
A complete proof would follow from any argument producing a finite set of primes $S$ such
that, for one direction and every $n\ge1$,
\[
 \supp(A^n-B^n)\setminus S
 \subseteq \supp(C^n-D^n).
\]
\end{criterion}

This target is qualitative and may be more accessible than the exponential gcd target of
Criterion~\ref{crit:cyclotomic-propagation}.  On the other hand, it demands simultaneous
control for every $n$.  Zsigmondy supplies new input primes, but the original hypotheses do
not say that the corresponding output order divides $n$.  The next section quantifies the
least multiplier needed to enforce that divisibility.

\section{Finite-field capture cost}
\label{sec:capture}

Support inclusion asks whether an input order kills the output.  When it does not, there is
still a least positive multiple of the input order that does.  This least multiplier is the
central new invariant of the report.

\subsection{Removing the finitely many nonunit primes}

Retain a direction $r=A/B>1$, $s=C/D=r^x>1$ in lowest terms.  For
$U_n=A^n-B^n$, define the $CD$-clean part
\[
 \widetilde U_n
 :=\prod_{\substack{p^{e}\parallel U_n\\p\nmid CD}}p^e.
\]
Thus $\widetilde U_n$ is the largest divisor of $U_n$ relatively prime to $CD$.

\begin{lemma}[Cleaning costs only polynomial mass]
\label{lem:cleaning}
One has
\[
 \log\widetilde U_n=n\log A+O_{A,B,C,D}(\log(n+1)).
\]
\end{lemma}

\begin{proof}
The full gap satisfies $\log U_n=n\log A+o(1)$.  The quotient
$U_n/\widetilde U_n$ is supported on the fixed finite set of primes dividing $CD$.
Lemma~\ref{lem:fixed-prime-valuation} bounds the valuation at each such prime by
$O(\log(n+1))$.
\end{proof}

The cleaning is indispensable.  If a prime divides $C$ or $D$, then $CD^{-1}$ is not a unit
modulo that prime power, so no multiplicative order is defined.  Lemma~\ref{lem:cleaning}
shows that deleting these places loses no exponential information.

\subsection{Definition and exact order formula}

\begin{definition}[Full capture cost]
\label{def:capture}
If $\widetilde U_n>1$, let
\[
 t_n:=\ord_{\widetilde U_n}(CD^{-1})
\]
where $D^{-1}$ is taken modulo $\widetilde U_n$.  Define
\[
 \kappa_{r,s}(n):=\frac{t_n}{\gcd(t_n,n)}.
\]
If $\widetilde U_n=1$, set $\kappa_{r,s}(n)=1$.
\end{definition}

\begin{proposition}[Exact least-multiplier interpretation]
\label{prop:capture-formula}
The integer $\kappa_{r,s}(n)$ is the least $k\ge1$ such that
\[
 \widetilde U_n\mid C^{kn}-D^{kn}.
\]
Moreover,
\[
 1\le\kappa_{r,s}(n)
 \le \lambda(\widetilde U_n)
 \le \varphi(\widetilde U_n)
 \le \widetilde U_n,
\]
where $\lambda$ is the Carmichael function.
\end{proposition}

\begin{proof}
The displayed divisibility is equivalent to
$(CD^{-1})^{kn}\equiv1\pmod{\widetilde U_n}$, hence to $t_n\mid kn$.  The least positive
$k$ with this property is $t_n/\gcd(t_n,n)$.  The bounds follow from the definition of the
multiplicative order and the standard inequalities $t_n\le\lambda\le\varphi\le\widetilde
U_n$.
\end{proof}

For a prime power $p^e\parallel\widetilde U_n$, put
\[
 \kappa_{p^e}(n)
 :=\frac{\ord_{p^e}(CD^{-1})}
         {\gcd(\ord_{p^e}(CD^{-1}),n)}.
\]
The Chinese remainder theorem gives the local factorization
\[
 \kappa_{r,s}(n)
 =\operatorname{lcm}_{p^e\parallel\widetilde U_n}\kappa_{p^e}(n).
\]
At a primitive divisor $p$ of $A^n-B^n$, one has $\ord_p(A/B)=n$; the local cost is exactly
the excess in $\ord_p(C/D)$ not already absorbed by $n$.

\begin{proposition}[Support and cost]
\label{prop:support-cost}
For a fixed $n$, $\kappa_{r,s}(n)=1$ exactly when every prime-power factor of
$A^n-B^n$ outside the fixed set $\supp(CD)$ is reproduced by $C^n-D^n$.  In particular,
if $\kappa_{r,s}(n)=1$ for infinitely many unbounded $n$, then $x\in\Z$.
\end{proposition}

\begin{proof}
The first statement is Proposition~\ref{prop:capture-formula} with $k=1$.  If it occurs
infinitely often, then the reduced denominator of $V_n/U_n$ is supported on the fixed set of
primes dividing $CD$.  Theorem~\ref{thm:criteria-menu}(a) applies.
\end{proof}

\subsection{The phase transition}

\begin{theorem}[Capture-cost dichotomy]
\label{thm:capture-dichotomy}
Let $x>0$ satisfy $2^x,3^x\in\Z$, and fix a nonzero direction.
\begin{enumerate}[label=\textup{(\roman*)}]
\item If $x\in\Z$, then $\kappa_{r,s}(n)=1$ for every $n\ge1$.
\item If $x\notin\Z$, then
\[
 \kappa_{r,s}(n)\longrightarrow\infty.
\]
Equivalently, for every fixed $K$, all sufficiently large $n$ require a multiplier larger
than $K$ to capture the cleaned input gap.
\end{enumerate}
\end{theorem}

\begin{proof}
If $x=m\in\Z_{>0}$, then $C=A^m$, $D=B^m$ and
$A^n-B^n\mid A^{mn}-B^{mn}$, so $k=1$ is admissible.

Suppose $x\notin\Z$ and the costs fail to diverge.  Then some $K$ bounds
$\kappa_{r,s}(n)$ on an infinite set.  By the pigeonhole principle, a fixed
$1\le k\le K$ occurs for infinitely many $n$.  For those $n$,
\[
 \widetilde U_n\mid C^{kn}-D^{kn},
\]
so
\[
 \log\gcd(A^n-B^n,C^{kn}-D^{kn})
 \ge\log\widetilde U_n
 =n\log A+O(\log n).
\]
But $A/B$ and $(C/D)^k$ are multiplicatively independent.  Indeed, a dependence would give
$r^u=s^{kv}=r^{xkv}$ for nonzero integers $u,v$, hence $x=u/(kv)\in\Q$, and then
Lemma~\ref{lem:rational} would make $x$ integral.  The homogeneous gcd estimate applied to
$A/B$ and $(C/D)^k$ says that the logarithm of the same gcd is $o(n)$, a contradiction.
\end{proof}

\begin{resultbox}
The integrality transition is therefore exact:
\[
 x\in\Z\quad\Longleftrightarrow\quad
 \kappa_{r,s}(n)\equiv1
 \quad\Longleftrightarrow\quad
 \liminf_{n\to\infty}\kappa_{r,s}(n)<\infty.
\]
This holds in every fixed rational direction.
\end{resultbox}

\subsection{Partial capture cost}

Full divisibility is stronger than needed.  For $0<\delta<\log A$, define, for all
sufficiently large $n$,
\[
 \kappa_{r,s;\delta}(n)
 :=\min\left\{k\ge1:
 \log\gcd(A^n-B^n,C^{kn}-D^{kn})\ge\delta n\right\}.
\]
Existence follows because the full cost $k=\kappa_{r,s}(n)$ captures
$\widetilde U_n$, whose logarithm is $n\log A+O(\log n)$.

\begin{theorem}[Positive-mass capture also diverges]
\label{thm:partial-capture}
If $x\notin\Z$, then for every fixed $0<\delta<\log A$,
\[
 \kappa_{r,s;\delta}(n)\longrightarrow\infty.
\]
If $x\in\Z$, then $\kappa_{r,s;\delta}(n)=1$ for all sufficiently large $n$.
\end{theorem}

\begin{proof}
The integral case is immediate from $U_n\mid V_n$ and
$\log U_n=n\log A+o(1)$.  In the nonintegral case, boundedness on an infinite sequence would
again produce one fixed $k$.  The homogeneous gcd estimate for $r$ and $s^k$ would then give
$o(n)$ common logarithmic mass, contradicting the defining lower bound $\delta n$.
\end{proof}

Thus even recovering one percent of the exponential input mass requires a multiplier tending
to infinity.  This is stronger than the qualitative support defect statement and weaker than
full denominator saturation; it is often the most flexible target for a future argument.

\subsection{Ray covariance and divisibility in the layer index}

\begin{proposition}[Covariance on a lattice ray]
\label{prop:ray-covariance}
Let $m\ge1$.  If the direction $m\mathbf a$ is represented by the powered ratios
$r^m=A^m/B^m$ and $s^m=C^m/D^m$, then
\[
 Q_{m\mathbf a}(n)=Q_{\mathbf a}(mn),
 \qquad
 \kappa_{r^m,s^m}(n)=\kappa_{r,s}(mn).
\]
\end{proposition}

\begin{proof}
Both statements reduce to the identities
$(A^m)^n-(B^m)^n=A^{mn}-B^{mn}$ and
$(C^m)^{kn}-(D^m)^{kn}=C^{kmn}-D^{kmn}$.  The set of primes removed in the cleaning step is
unchanged by taking powers.
\end{proof}

\begin{proposition}[A weak layer-divisibility inequality]
\label{prop:layer-inequality}
If $m\mid n$ and $q=n/m$, then
\[
 \kappa_{r,s}(m)\le q\,\kappa_{r,s}(n).
\]
\end{proposition}

\begin{proof}
The cleaned divisor $\widetilde U_m$ divides $\widetilde U_n$.  If
$k=\kappa_{r,s}(n)$, then $\widetilde U_n$ divides $C^{kn}-D^{kn}$, so
$\widetilde U_m$ divides $C^{kqm}-D^{kqm}$.  Therefore $qk$ is admissible at level $m$.
\end{proof}

The inequality is too weak to force boundedness, but it is useful in computations and in a
possible compactness argument on highly divisible indices.

\subsection{Simultaneous synchronization of the two bases}

Apply the construction to the diagonal directions
\[
 (r,s)=(2,M),\qquad (r,s)=(3,N),
\]
and write $\kappa_2(n)$ and $\kappa_3(n)$.  Set
\[
 \kappa_{\mathrm{syn}}(n):=\operatorname{lcm}(\kappa_2(n),\kappa_3(n)).
\]

\begin{theorem}[Two-base synchronization phase transition]
\label{thm:synchronization}
For every positive solution $x$:
\begin{enumerate}[label=\textup{(\roman*)}]
\item if $x\in\Z$, then $\kappa_{\mathrm{syn}}(n)=1$ for all $n$;
\item if $x\notin\Z$, then $\kappa_{\mathrm{syn}}(n)\to\infty$;
\item for every $n$,
\[
 \kappa_{\mathrm{syn}}(n)<6^n;
\]
\item with $L=\kappa_{\mathrm{syn}}(n)$, the same multiplier simultaneously gives
\begin{align*}
 \widetilde{(2^n-1)}^{\,(M)}&\mid M^{Ln}-1,\\
 \widetilde{(3^n-1)}^{\,(N)}&\mid N^{Ln}-1,
\end{align*}
where the tildes remove prime powers dividing $M$ and $N$, respectively, and their logarithms
are $n\log2+O(\log n)$ and $n\log3+O(\log n)$.
\end{enumerate}
\end{theorem}

\begin{proof}
Parts (i) and (ii) follow from Theorem~\ref{thm:capture-dichotomy}.  The individual bounds
$\kappa_2(n)\le2^n-1$ and $\kappa_3(n)\le3^n-1$ imply
\[
 \kappa_{\mathrm{syn}}(n)
 \le\kappa_2(n)\kappa_3(n)<6^n.
\]
Because $L$ is a multiple of each individual least cost, both divisibilities in (iv) hold.
The logarithmic sizes follow from Lemma~\ref{lem:cleaning}.
\end{proof}

\begin{criterion}[Archimedean-to-Frobenius bridge]
\label{crit:bridge}
The \AEconj{} conjecture would follow if the simultaneous real relation
\[
 \log M=x\log2,\qquad \log N=x\log3
\]
could be used to prove that $\kappa_{\mathrm{syn}}(n)$ is bounded on an unbounded sequence of
$n$.  More generally, it is enough to bound both partial costs
$\kappa_{2,M;\delta_2}(n)$ and $\kappa_{3,N;\delta_3}(n)$ for any fixed positive
$\delta_2,\delta_3$.
\end{criterion}

This is the most concentrated open bridge found in the present work.  The hypotheses are
archimedean identities among logarithms; the conclusion concerns orders in finite
multiplicative groups.  Existing transcendence estimates control how closely integer powers
can approach at the real place, but they do not directly bound these finite-field order
multipliers.

\subsection{Collapse to the primitive core}

\begin{proposition}[Base-adic depth is invisible to capture]
\label{prop:capture-core}
If $M=2^aR$ with $a\ge0$, then the diagonal capture cost for $(2,M)$ equals that for
$(2,R)$ at every layer.  If $N=3^bS$, then the cost for $(3,N)$ equals that for $(3,S)$.
\end{proposition}

\begin{proof}
Modulo $2^n-1$ one has
\[
 M^{kn}=2^{akn}R^{kn}\equiv R^{kn}.
\]
Also, because $2\nmid2^n-1$, removing the prime factors of $M$ removes exactly the prime
factors of $R$ that occur in $2^n-1$.  Hence the cleaned modulus and the admissible values of
$k$ are the same.  The base-$3$ argument is identical.
\end{proof}

For a least hypothetical counterexample in the primitive normalization of the unified
report, the synchronization invariant therefore sees only the odd/base-free cores.  The
large powers of $2$ and $3$ attached to those cores are arithmetically invisible at their
own cyclotomic moduli.

\section{Exact numerical reconnaissance}
\label{sec:computation}

The theorems above are asymptotic and ineffective: the Subspace Theorem does not provide a
practical threshold beyond which a chosen cost exceeds a chosen bound.  Exact computation is
therefore useful for calibrating the invariant, but it cannot certify the asymptotic regime.
This section uses small independent pairs as surrogates and an integral pair as a control.

\subsection{Control and surrogate pairs}

For the input base $2$ consider three outputs:
\[
 M=4,\qquad M=3,\qquad M=5.
\]
The first is the integral case $4=2^2$, so $\kappa_{2,4}(n)=1$ identically.  The other two
are multiplicatively independent from $2$ and therefore model the finite-field behavior that
a noninteger solution would be forced to have in the diagonal direction.  They are not
claimed to extend to solutions of the simultaneous two-base problem.

\begin{experiment}[Selected full-capture costs]
\label{exp:cost-table}
Exact factorization of $2^n-1$ and exact multiplicative orders give:
\[
\begin{array}{c|r|r|r}
\toprule
n & \kappa_{2,4}(n) & \kappa_{2,3}(n) & \kappa_{2,5}(n)\\
\midrule
5  &1&6&3\\
7  &1&18&6\\
11 &1&8&4\\
13 &1&70&105\\
17 &1&7710&3855\\
32 &1&2048&2048\\
40 &1&1542&1542\\
60 &1&110&55\\
\bottomrule
\end{array}
\]
\end{experiment}

The independent costs are emphatically not monotone.  For example,
$\kappa_{2,3}(17)=7710$ but $\kappa_{2,3}(60)=110$.  Theorem~\ref{thm:capture-dichotomy}
says only that every fixed bound is eventually left behind; it does not assert monotonicity,
and the ineffective threshold may be very large.

\subsection{Local explanations for several rows}

At prime indices with $2^n-1$ prime, the cost is especially transparent:
\begin{align*}
 n=5:&\quad 2^5-1=31,
 &\ord_{31}(3)=30,&\quad \kappa_{2,3}(5)=30/5=6;\\
 n=7:&\quad 2^7-1=127,
 &\ord_{127}(3)=126,&\quad \kappa_{2,3}(7)=126/7=18;\\
 n=13:&\quad 2^{13}-1=8191,
 &\ord_{8191}(3)=910,&\quad \kappa_{2,3}(13)=910/13=70;\\
 n=17:&\quad 2^{17}-1=131071,
 &\ord_{131071}(3)=131070,&\quad \kappa_{2,3}(17)=7710.
\end{align*}
For $M=4$, by contrast,
$\ord_{2^n-1}(4)$ always divides $n$, because $4^n=2^{2n}\equiv1$; the cost is $1$.

Composite layers show the lcm nature of the invariant.  At $n=32$,
\[
 2^{32}-1=3\cdot5\cdot17\cdot257\cdot65537.
\]
After deleting the factor $3$ for the output $M=3$, the order of $3$ modulo the cleaned
modulus is $65536$, and hence
\[
 \kappa_{2,3}(32)=\frac{65536}{\gcd(65536,32)}=2048.
\]
At $n=60$, many local orders are already absorbed by the highly divisible layer index, which
explains the much smaller cost $110$.

\subsection{Denominator ratios at finite height}

For an output $M$ define
\[
 \rho_M(n)
 :=\frac{\log\bigl((2^n-1)/\gcd(2^n-1,M^n-1)\bigr)}{n\log2}.
\]
For multiplicatively independent $2$ and $M$, the BCZ theorem gives
$\rho_M(n)\to1$, but convergence can be irregular.  Some exact values, rounded only after
the exact gcd computation, are:
\[
\begin{array}{c|cc|cc}
\toprule
n & \gcd(2^n-1,3^n-1)&\rho_3(n)
  & \gcd(2^n-1,5^n-1)&\rho_5(n)\\
\midrule
11&23&0.588703&1&0.999936\\
17&1&0.999999&1&0.999999\\
32&85&0.799707&51&0.822737\\
40&11275&0.663479&1353&0.739951\\
60&47322275&0.575066&698476779&0.510339\\
\bottomrule
\end{array}
\]
The dips do not contradict the limit.  Subexponential gcds can still be very large at
moderate indices, and known lower-bound constructions show that large exceptional gcds occur
infinitely often for independent bases.  This warns against inferring asymptotics from a
short scan.

\subsection{What should be computed for an actual candidate}

A hypothetical solution would supply the exact integers $M$ and $N$.  The most informative
finite experiment would then be:
\begin{enumerate}[label=\textup{(\roman*)}]
\item compute $\kappa_2(n)$ and $\kappa_3(n)$ on prime, prime-power, and highly composite
indices;
\item factor the local costs into excess order primes;
\item compute the partial costs for several fractions
$\delta/\log2$ and $\delta/\log3$;
\item search for directions $(a,b)$ of small height with anomalously small
$\kappa_{\mathbf a}(n)$;
\item test whether the same small multiplier recurs across both bases or across several
adjacent smooth directions.
\end{enumerate}
A bounded recurring multiplier would not merely be suggestive; by
Theorem~\ref{thm:capture-dichotomy} it would be a rigorous contradiction certificate once an
unbounded family of layers were proved.

\subsection{Reproducible exact code}

The following code generated Table~\ref{exp:cost-table}.  It uses no floating-point
arithmetic in the factorization, order, gcd, or cost calculations.

\begin{lstlisting}[style=pythonstyle,caption={Exact full-capture cost for a diagonal pair.},label={lst:capture}]
from __future__ import annotations

from math import gcd
from typing import Dict, Tuple

import sympy as sp


def cleaned_gap(base: int, output: int, n: int) -> Tuple[int, Dict[int, int]]:
    """Return the part of base**n - 1 coprime to output, with its factorization."""
    if base < 2 or output < 2 or n < 1:
        raise ValueError("base, output >= 2 and n >= 1 are required")

    gap = base**n - 1
    factors = sp.factorint(gap)
    modulus = 1
    kept: Dict[int, int] = {}
    for p, e in factors.items():
        if output % p != 0:
            modulus *= p**e
            kept[int(p)] = int(e)
    return modulus, kept


def full_capture_cost(base: int, output: int, n: int) -> int:
    """Least k >= 1 for which the cleaned gap divides output**(k*n) - 1."""
    modulus, _ = cleaned_gap(base, output, n)
    if modulus == 1:
        return 1

    order = int(sp.n_order(output % modulus, modulus))
    return order // gcd(order, n)


def verify_cost(base: int, output: int, n: int, k: int) -> bool:
    """Certify admissibility and minimality of a proposed full capture cost."""
    modulus, _ = cleaned_gap(base, output, n)
    if pow(output, k * n, modulus) != 1 % modulus:
        return False
    return all(pow(output, j * n, modulus) != 1 % modulus for j in range(1, k))


if __name__ == "__main__":
    indices = [5, 7, 11, 13, 17, 32, 40, 60]
    for n in indices:
        row = [n]
        for output in (4, 3, 5):
            k = full_capture_cost(2, output, n)
            assert verify_cost(2, output, n, k)
            row.append(k)
        print(*row)
\end{lstlisting}

For large $n$, complete factorization of $A^n-B^n$ becomes the bottleneck.  The local formula
permits partial certified lower bounds: factor any divisor $F\mid\widetilde U_n$, compute the
order modulo each known prime power, and take the lcm of the resulting local costs.  This
lower bound is enough to rule out small synchronization multipliers without factoring the
entire gap.

\section{Hadamard preconditioning: what it can and cannot change}
\label{sec:hadamard}

The prompt highlights recent Hadamard-matrix computations, so it is natural to ask whether a
Hadamard sign transform could strengthen one of the determinant arguments in the unified
report.  There is a clean answer at the global level.

\subsection{Analytic neutrality under orthogonal normalization}

Let $H$ be a Hadamard matrix of order $m$:
\[
 H\in\{\pm1\}^{m\times m},\qquad HH^{\mathsf T}=mI_m.
\]
Then $U=m^{-1/2}H$ is orthogonal.

\begin{theorem}[Exterior-power neutrality]
\label{thm:hadamard-neutrality}
Let $A$ be a real $m\times n$ matrix and $1\le k\le\min(m,n)$.  Then:
\begin{enumerate}[label=\textup{(\roman*)}]
\item $UA$ and $A$ have the same singular values;
\item
\[
 \|\wedge^k(UA)\|_{\mathrm F}=\|\wedge^k A\|_{\mathrm F};
\]
\item equivalently,
\[
 \sum_{\substack{I\subseteq[m],\ |I|=k\\J\subseteq[n],\ |J|=k}}
 \det((HA)_{I,J})^2
 =m^k
 \sum_{\substack{I\subseteq[m],\ |I|=k\\J\subseteq[n],\ |J|=k}}
 \det(A_{I,J})^2.
\]
If $A$ is square, then
\[
 |\det(HA)|=m^{m/2}|\det A|.
\]
\end{enumerate}
\end{theorem}

\begin{proof}
Left multiplication by the orthogonal matrix $U$ preserves $A^{\mathsf T}A$, hence all
singular values.  The singular values of $\wedge^k A$ are the $k$-fold products of singular
values of $A$, proving (ii).  The squared Frobenius norm of $\wedge^k A$ is the sum of the
squares of all $k\times k$ minors, giving (iii) after $H=\sqrt m\,U$.  Finally,
$|\det H|=m^{m/2}$ follows from $HH^{\mathsf T}=mI$.
\end{proof}

Thus Hadamard mixing can redistribute determinant mass among individual minors, but it cannot
increase the aggregate exterior-power budget beyond the explicit factor $m^{k/2}$ already
paid for by the row scaling.  Any analytic determinant estimate based only on singular
values, Frobenius norms, or the sum of squared minors is invariant after normalization.

\subsection{$p$-adic neutrality away from the order}

There is an equally useful arithmetic statement.

\begin{theorem}[Local Smith-profile neutrality]
\label{thm:hadamard-smith}
Let $A$ be an integer $m\times n$ matrix and $H$ a Hadamard matrix of order $m$.  For every
prime $p\nmid m$, left multiplication by $H$ preserves the Smith invariants of $A$ over
$\Z_p$.  Equivalently, for every $k$, the $p$-adic valuation of the determinantal ideal
generated by all $k\times k$ minors is the same for $A$ and $HA$.
\end{theorem}

\begin{proof}
Since
\[
 H^{-1}=\frac1m H^{\mathsf T},
\]
the matrix $H$ lies in $\mathrm{GL}_m(\Z_p)$ whenever $p\nmid m$.  Multiplication by an
invertible matrix over $\Z_p$ is a sequence of invertible row operations, which preserves
the local Smith normal form and every Fitting/determinantal ideal.
\end{proof}

\begin{corollary}[Only primes dividing the order can be globally gained]
Any improvement in the common $p$-adic divisibility of all $k$-minors after Hadamard
preconditioning must occur at a prime $p\mid m$.  At primes $p\nmid m$, only the distribution
among individual minors can change; the minimum valuation and the generated ideal cannot.
\end{corollary}

This sharply limits the role of a newly available Hadamard order.  Choosing $m$ with many
prime factors can create local room at those primes, but it also introduces the known global
determinant factor $m^{m/2}$.  A proof must exploit a specially selected minor or congruence,
not generic orthogonal flattening.

\subsection{A viable selective-minor program}

The neutrality theorem does not make Hadamard matrices useless.  It identifies the only
remaining mechanism:
\begin{enumerate}[label=\textup{(\roman*)}]
\item form a structured interpolation or moment matrix $A$ from the integer values
$2^x,3^x$ and their semigroup powers;
\item choose Hadamard rows so that one designated minor experiences systematic cancellation
at the real place;
\item retain enough integrality that the same minor is divisible by a large known integer;
\item prove that this minor is nonzero and smaller in absolute value than its divisibility
modulus.
\end{enumerate}
The exterior-power identity says that any exceptionally small chosen minor is balanced by
larger minors elsewhere.  That is acceptable: a determinant contradiction needs one minor,
not an average.  The difficult step is to align the small analytic minor with a strong
arithmetic divisor.  No such alignment has yet been found.

\begin{criterion}[Selective Hadamard certificate]
\label{crit:hadamard}
A Hadamard-based proof must use information not invariant under left orthogonal action:
a designated row subset, a sign-sensitive congruence, or a prime dividing the Hadamard
order.  Arguments depending only on the full determinant, singular values, or aggregate
minor norms cannot improve the existing determinant balance.
\end{criterion}

\section{Attempts to close the synchronization bridge}
\label{sec:attempts}

The preceding sections isolate several exact proof certificates.  This section records the
main attempts made to derive one of them from the simultaneous integrality hypotheses, and
pinpoints why each attempt presently stops.  These are not excuses; they are constraints that
a successful continuation must explicitly overcome.

\subsection{Finite packets of directions}

For $H\ge1$, let
\[
 \mathcal F_H
 :=\{(a,b)\in\Z^2\setminus\{0\}:|a|+|b|\le H\}/\{\pm1\},
\]
with each sign class oriented so that $2^a3^b>1$.

\begin{proposition}[Uniform escape on every fixed packet]
\label{prop:finite-packet}
If $x$ is a hypothetical noninteger solution, then for each fixed $H$,
\[
 \min_{\mathbf a\in\mathcal F_H}\kappa_{\mathbf a}(n)\longrightarrow\infty.
\]
The same assertion holds for every fixed positive-mass partial cost.
\end{proposition}

\begin{proof}
The packet is finite, and the cost in each fixed direction tends to infinity by
Theorems~\ref{thm:capture-dichotomy} and~\ref{thm:partial-capture}.  Given $K$, take the
maximum of the finitely many thresholds beyond which each cost exceeds $K$.
\end{proof}

Thus a counterexample cannot hide a permanently cheap small direction.  A direction chosen
adaptively with $n$ could still have low cost, and the current Subspace-Theorem argument is
not uniform when the height of that direction grows.  One possible continuation is therefore
a quantitative, height-uniform version of the directional gcd theorem.

\begin{question}[Growing-height uniformity]
Can one prove a function $H(n)\to\infty$ for which a noninteger solution would force
\[
 \min_{\mathbf a\in\mathcal F_{H(n)}}\kappa_{\mathbf a}(n)\to\infty?
\]
A sufficiently rapid $H(n)$ could be confronted with geometry-of-numbers constructions of
short lattice directions.
\end{question}

\subsection{The three-gap triangle}

The adjacent smooth nodes $1,2,3,4$ produce three input gaps
\[
 U_0=2^n-1,
 \qquad U_1=3^n-2^n,
 \qquad U_2=4^n-3^n.
\]
They satisfy
\[
 U_0+U_1+U_2=4^n-1,
 \qquad U_1+U_2=2^nU_0.
\]
Their output analogues are
\[
 V_0=M^n-1,
 \qquad V_1=N^n-M^n,
 \qquad V_2=M^{2n}-N^n,
\]
with
\[
 V_0+V_1+V_2=M^{2n}-1,
 \qquad V_1+V_2=M^nV_0.
\]

\begin{proposition}[Exact pairwise overlap of the input triangle]
\label{prop:triangle-gcd}
For every $n\ge1$,
\[
 \gcd(U_0,U_1)=\gcd(U_0,U_2)=\gcd(U_1,U_2)
 =\gcd(2^n-1,3^n-1).
\]
Consequently the logarithm of every pairwise gcd is $o(n)$.
\end{proposition}

\begin{proof}
Modulo $U_0$ one has $2^n\equiv4^n\equiv1$, so both $U_1$ and $-U_2$ are congruent to
$3^n-1$.  This proves the first two identities.  Modulo $U_1$ one has
$3^n\equiv2^n$, hence
\[
 U_2\equiv4^n-2^n=2^n(2^n-1).
\]
The odd integer $U_1$ is coprime to $2^n$, so
$\gcd(U_1,U_2)=\gcd(U_1,2^n-1)$, which is the first gcd.  The final estimate is BCZ for the
independent bases $2$ and $3$.
\end{proof}

\begin{corollary}[An almost-coprime gap packet]
\label{cor:triangle-lcm}
Let $L_n=\operatorname{lcm}(U_0,U_1,U_2)$.  Then
\[
 \log L_n=n\log24+o(n).
\]
\end{corollary}

\begin{proof}
The three gaps have logarithms $n\log2+o(n)$, $n\log3+o(n)$, and
$n\log4+o(n)$.  Inclusion--exclusion for valuations, together with the $o(n)$ pairwise
overlaps, gives the result; the triple overlap is bounded by a pairwise overlap.
\end{proof}

This is promising because the adjacent gaps jointly carry almost $24^n$ independent modulus.
However, the telescoping output identities do not force any one input gap to divide its
output counterpart.  Writing $Q_i=V_i/U_i$ gives
\[
 Q_1U_1+Q_2U_2=M^nQ_0U_0,
\]
but each product $Q_iU_i=V_i$ is already integral.  The large, almost-coprime denominators
are multiplied by precisely the gaps that cancel them, so the identity supplies no new
cross-divisibility.  A successful triangle argument needs an additional relation among the
$Q_i$ with coefficients that do not contain the full $U_i$.

\subsection{Real near-collisions do not control finite-field orders}

Take integers $u,v>0$ and put
\[
 \varepsilon=u\log2-v\log3.
\]
Then
\begin{align*}
 2^u-3^v&=3^v(\ee^{\varepsilon}-1),\\
 M^u-N^v&=N^v(\ee^{x\varepsilon}-1).
\end{align*}
Good rational approximations to $\log_2 3$ make both relative gaps small, and the unified
report extensively develops the resulting secant and determinant information.  For the
capture problem, however, one needs congruences such as
\[
 M^{kn}\equiv1\pmod{2^n-1}.
\]
The size of $M^u-N^v$ at the real place does not control its residue modulo a prime divisor of
$2^n-1$.  Even an exponentially small relative error may correspond to an enormous nonzero
integer difference containing none of the required primes.

There is one rigorous way a near-collision could help: prove that a product of desired input
prime powers exceeds the absolute value of a nonzero integer known to be divisible by that
product.  Existing approximations do not achieve this balance.  Their integer differences
have exponential scale comparable to the moduli one wants to capture, and the known lower
bounds for linear forms in logarithms prevent the extreme cancellation needed.

\subsection{Why simultaneous bases have not yet bounded the cost}

The two diagonal costs are built from
\[
 \ord_{p^e}(M)\quad(p^e\mid2^n-1),
 \qquad
 \ord_{q^f}(N)\quad(q^f\mid3^n-1).
\]
The identity
\[
 \frac{\log M}{\log2}=\frac{\log N}{\log3}=x
\]
links these integers only through the real embedding.  There is no operation ``raising to
$x$'' in the finite cyclic groups unless $x$ is already known to be integral.  In particular,
$2^n\equiv1\pmod p$ does not imply any formal congruence for $M^n$.

A tempting Chinese-remainder construction chooses
$L=\kappa_{\mathrm{syn}}(n)$, for which nearly all of both input gaps are captured at exponent
$Ln$.  This merely proves the existence of an $L<6^n$; Euler's theorem already supplies such
an exponential bound.  The gcd estimate rules out bounded $L$, but not a growing $L$.
Uniform estimates for
\[
 \gcd(2^n-1,M^{Ln}-1)
\]
with both exponents varying would have to be strong enough to compare an $e^{\Theta(n)}$ gcd
against a second exponent $Ln$.  The qualitative Subspace Theorem estimate loses control as
$L$ grows.

\subsection{Recurrence extraction: the exact obstruction}

The gcd sequence $g_{2,M}(n)$ is automatically a divisibility sequence.  One might hope that
the simultaneous presence of $N=3^x$ forces it to satisfy a recurrence.  Theorem~\ref{thm:ND}
shows that this hope, if realized, would immediately prove the conjecture; it also shows that
there is no weaker nonperiodic LDS hidden inside the diagonal gcd under a counterexample.

The missing implication is therefore stark:
\[
 2^x,3^x\in\Z
 \quad\stackrel{?}{\Longrightarrow}\quad
 g_{2,2^x}(n)\text{ has a nonperiodic recurrent factor.}
\]
The local support of a gcd sequence is governed by multiplicative orders and is typically an
infinite union of arithmetic progressions.  A constant-coefficient recurrence has only
finitely many characteristic roots.  No mechanism is presently known that compresses the
order data supplied by the second base into such a finite spectral set.

\subsection{Primitive divisors: freshness without size}

For every large layer, Zsigmondy provides a prime with new input order.  At such a prime $p$,
the local cost is
\[
 \frac{\ord_p(s)}{\gcd(\ord_p(s),n)}.
\]
Corollary~\ref{cor:defect-infinite} ensures infinitely many instances where this exceeds $1$;
Theorem~\ref{thm:capture-dichotomy} forces unbounded aggregate excess.  What is missing is a
lower bound for the size or abundance of primitive primes with large excess output order.
The standard primitive-divisor theorem is qualitative and allows one fresh prime to carry
negligible logarithmic mass.

A potentially decisive strengthening would be:
\begin{quote}
For one fixed direction, a positive proportion of
$\log\Phi_n(A,B)$ is carried by prime powers $p^e$ for which
$\ord_{p^e}(C/D)\nmid Kn$ for every fixed $K$.
\end{quote}
The current gcd theorem proves the aggregate conclusion for each fixed $K$, but only after
summing over all layers; it does not tie it to a structure coming from the four original
integer values.

\subsection{Most promising next technical target}

Among the available bridges, the following appears both sharp and narrowly formulated.

\begin{conjecture}[Bounded synchronized positive-mass capture]
\label{conj:bounded-capture}
If $M=2^x$ and $N=3^x$ are integers, then there exist
$\delta_2,\delta_3>0$, a fixed $K$, and infinitely many $n$ for which one can find
$1\le k\le K$ satisfying
\begin{align*}
 \log\gcd(2^n-1,M^{kn}-1)&\ge\delta_2n,\\
 \log\gcd(3^n-1,N^{kn}-1)&\ge\delta_3n.
\end{align*}
\end{conjecture}

The conjecture is automatic with $k=1$ when $x$ is integral and impossible for a noninteger
solution by Theorem~\ref{thm:partial-capture}.  Proving it would settle \AEconj{}.  Its advantage
over the original statement is that it asks for a concrete finite-field phenomenon that can
be attacked using multiplicative orders, Chebotarev-type distribution, or structured
resultants.

\section{Lean formalization blueprint}
\label{sec:lean}

A formal development should not begin with the full transcendence statement.  Most of the
new work is an elementary finite arithmetic package that can be verified independently of
the deep gcd theorems.  This section gives a dependency-aware plan designed to expose the
precise trusted boundary.

\subsection{Proposed module structure}

\begin{longtable}{p{0.25\textwidth}p{0.67\textwidth}}
\toprule
Module & Content \\
\midrule
\endfirsthead
\toprule
Module & Content \\
\midrule
\endhead
\texttt{GapDirection} & Reduced positive rational directions; homogeneous gaps
$A^n-B^n$; powering a direction; covariance along a lattice ray. \\
\texttt{FixedPrimeValuation} & LTE-based bound
$v_p(A^n-B^n)\le C+v_p(n)$, including the two-adic branch. \\
\texttt{CleanedGap} & Largest divisor of a gap coprime to a fixed integer; factorization and
logarithmic-loss lemmas. \\
\texttt{CaptureCost} & Multiplicative order modulo the cleaned gap; least-multiplier formula;
local lcm decomposition; elementary upper bounds. \\
\texttt{SupportDefect} & Equivalence between support inclusion and divisibility of finite-field
orders; forward and reverse defect sets. \\
\texttt{GapTriangle} & Exact pairwise gcd identities for the nodes $1,2,3,4$ and the lcm
consequence, conditional on an abstract subexponential-gcd input. \\
\texttt{HadamardPrecondition} & Orthogonal normalization, determinant scaling, exterior-power
norm identity, and local Smith invariance for primes not dividing the order. \\
\texttt{AEBridge} & Construction of $A,B,C,D$ from $(a,b)$ and a positive solution
$M=2^x,N=3^x$; rationality detection from multiplicative dependence. \\
\texttt{ExternalGCD} & Clearly marked theorem interfaces for Corvaja--Zannier and
Nguyen-Dang. \\
\texttt{Consequences} & Denominator saturation, capture divergence, fixed-support criteria,
and recurrence signatures derived from the interfaces. \\
\bottomrule
\end{longtable}

\subsection{Core data structure}

One convenient abstraction is:
\begin{lstlisting}[style=pythonstyle,language={},caption={Lean-style pseudocode for a reduced direction.}]
structure RatDirection where
  A B C D : Nat
  hB : 0 < B
  hD : 0 < D
  hAB : Nat.Coprime A B
  hCD : Nat.Coprime C D
  hABpos : B < A
  hCDpos : D < C

namespace RatDirection

def inputGap  (R : RatDirection) (n : Nat) : Nat := R.A^n - R.B^n
def outputGap (R : RatDirection) (n : Nat) : Nat := R.C^n - R.D^n

def cleanedGap (R : RatDirection) (n : Nat) : Nat :=
  (R.inputGap n).coprimePart (R.C * R.D)

-- Defined through the order of C * D^{-1} in ZMod (cleanedGap n).
def captureCost (R : RatDirection) (n : Nat) : Nat := ...

end RatDirection
\end{lstlisting}
The actual implementation may represent the unit $CD^{-1}$ in
$(\mathbb Z/\widetilde U_n\mathbb Z)^\times$ rather than in the ambient ring, avoiding
special cases in the order API.

\subsection{Elementary theorem milestones}

The first formal milestone should require no unproved number-theoretic interface:

\begin{enumerate}[label=\textbf{L\arabic*.}]
\item \textbf{Least multiplier.}  Prove
\[
 \kappa(n)=\min\{k\ge1:\widetilde U_n\mid C^{kn}-D^{kn}\}
 =\frac{t_n}{\gcd(t_n,n)}.
\]
\item \textbf{Local decomposition.}  Prove that the global cost is the lcm of its prime-power
costs.
\item \textbf{Integral branch.}  If $C=A^m,D=B^m$, prove $\kappa(n)=1$.
\item \textbf{Ray covariance.}  Prove
$\kappa_{r^m,s^m}(n)=\kappa_{r,s}(mn)$.
\item \textbf{Layer inequality.}  For $m\mid n$, prove
$\kappa(m)\le(n/m)\kappa(n)$.
\item \textbf{Core collapse.}  Formalize the congruences
$(2^aR)^{kn}\equiv R^{kn}\pmod{2^n-1}$ and their base-$3$ analogues.
\item \textbf{Gap triangle.}  Prove all three pairwise gcd identities by modular rewriting.
\item \textbf{Hadamard local neutrality.}  Over $\mathbb Z_p$, construct the inverse
$m^{-1}H^{\mathsf T}$ for $p\nmid m$ and transfer determinantal ideals.
\end{enumerate}

These statements are exact and computationally testable.  Completing them would already
produce a reusable finite-field library independent of the status of \AEconj{}.

\subsection{Deep theorem interfaces}

The formal proof should isolate, not silently assume, three external results.

\begin{enumerate}[label=\textbf{E\arabic*.}]
\item \textbf{Homogeneous gcd estimate.}  For multiplicatively independent reduced positive
rationals $A/B,C/D$,
\[
 \log\gcd(A^n-B^n,C^n-D^n)=o(n).
\]
A first Lean development can expose this as an axiomatically imported theorem with a precise
bibliographic comment.  A full formalization would require a major formal Subspace-Theorem
project.
\item \textbf{Support problem.}  Almost-everywhere order divisibility implies an integral
power relation.  Formalizing the published proof would require algebraic number theory and
Chebotarev machinery well beyond the elementary core.
\item \textbf{Recurrence classification.}  Nguyen-Dang's 2026 theorem should be a separate
interface because it additionally depends on modern classification of integer linear
divisibility sequences.
\end{enumerate}

No final theorem should be labeled ``fully formalized'' while any of these interfaces remains
an unverified axiom.  The elementary consequences can nevertheless be machine-checked
relative to clearly named assumptions.

\subsection{Formal statement of the capture dichotomy}

A practical theorem interface could look schematically as follows:
\begin{lstlisting}[style=pythonstyle,language={},caption={Lean-style pseudocode for the main conditional theorem.}]
axiom homogeneous_gcd_isLittleO
  (R : RatDirection)
  (hindep : MultiplicativelyIndependent (R.A / R.B) (R.C / R.D)) :
  IsLittleO atTop
    (fun n => Real.log (Nat.gcd (R.inputGap n) (R.outputGap n)))
    (fun n => (n : Real))

theorem captureCost_tendsto_atTop
  (R : RatDirection)
  (hindep : MultiplicativelyIndependent (R.A / R.B) (R.C / R.D)) :
  Tendsto R.captureCost atTop atTop := by
  -- contradiction: a bounded subsequence gives a fixed multiplier k;
  -- the cleaned input gap then lies in a forbidden exponential gcd.
  ...
\end{lstlisting}
The mathematical proof uses an infinite pigeonhole step.  In Lean it may be cleaner to prove
the contrapositive: for each fixed $K$, there exists $N$ such that no $1\le k\le K$ captures
the cleaned gap for $n\ge N$, then take the maximum of the finitely many thresholds attached
to $k=1,\dots,K$.

\subsection{Certified computation}

For concrete inputs, a certificate need not trust a factorization oracle.  It can include:
\begin{enumerate}[label=\textup{(\roman*)}]
\item a prime-power factorization of a divisor $F\mid A^n-B^n$;
\item primality certificates for the listed primes;
\item modular exponentiation showing the claimed local order;
\item for each proper prime divisor of that order, a nonidentity power proving minimality.
\end{enumerate}
The lcm of the certified local costs is then a formally verified lower bound for
$\kappa(n)$.  This is enough to establish that no multiplier up to $K$ works, without a full
factorization of the input gap.

\subsection{What a final Lean theorem would still need}

After all modules above, the ultimate statement would have the form
\begin{lstlisting}[style=pythonstyle,language={}]
theorem alaoglu_erdos
  (x : Real)
  (h2 : IsInteger (Real.rpow 2 x))
  (h3 : IsInteger (Real.rpow 3 x)) :
  IsInteger x := by
  ...
\end{lstlisting}
The unresolved line remains the derivation of a bounded recurrence/support/capture
certificate from \texttt{h2} and \texttt{h3}.  Formalization will make this gap impossible to
blur, which is a feature rather than a defect.

\section{Conclusions and prioritized continuation plan}
\label{sec:conclusion}

No complete proof has been obtained.  The main progress is a change of coordinates that
turns the conjecture into an exact finite-place phase transition.

\begin{resultbox}
For every fixed rational lattice direction, integer exponents make every power-gap quotient
integral and every capture cost equal to $1$.  A hypothetical noninteger solution forces the
opposite extreme: the quotient denominator retains asymptotically all input exponential
mass, support defects occur at infinitely many primes in both orientations, every fixed
multiplier captures only $o(n)$ logarithmic mass, and the least full or positive-mass
multiplier tends to infinity.
\end{resultbox}

This package gives several independent ways to finish:
\begin{enumerate}[label=\textbf{P\arabic*.}]
\item \textbf{Bounded positive-mass synchronization.}  Prove
Conjecture~\ref{conj:bounded-capture}.  This is the highest-priority target because it asks
only for a fixed positive fraction of two cyclotomic gaps, not full divisibility.
\item \textbf{Almost-everywhere support inclusion.}  Find a finite set $S$ for which one
diagonal support is contained in its output support outside $S$.  The support theorem then
finishes immediately.
\item \textbf{A nonperiodic common LDS.}  Use both integer outputs to construct any
nonperiodic recurrent factor of $2^n-1$ and $M^n-1$.  Nguyen-Dang's theorem turns this into a
short certificate.
\item \textbf{Uniformity in growing directions.}  Quantify the homogeneous gcd estimate as
the height of $(a,b)$ grows, then combine it with a geometry-of-numbers choice of a short
direction having unusually low capture cost.
\item \textbf{Selective Hadamard minor.}  Search only for sign patterns that align one
arithmetically divisible minor with analytic cancellation.  Do not spend effort on
orthogonally invariant determinant norms, which are neutral by
Theorem~\ref{thm:hadamard-neutrality}.
\end{enumerate}

The decisive conceptual issue is no longer hidden inside a large interpolation determinant.
It is the absence of a theorem transmitting a common real logarithmic scale to bounded
multiplicative-order distortion in finite fields.  Any claimed breakthrough should be
checked against this bridge: it must produce recurrence, support inclusion, exponential gcd
mass, or bounded capture in at least one fixed direction.

\begin{statusbox}{Final status on 22 August 2026: substantial new conditional rigidity and
several exact proof certificates; no proof of the \AEconj{} conjecture.}
\end{statusbox}

\appendix

\section{Nonduplication map relative to the unified report}
\label{app:nonduplication}

The attached source is a long unified manuscript incorporating twenty-seven continuations.
The present report was deliberately built around terminology and constructions absent from
that source.  A literal search of the supplied \texttt{.tex} found no occurrence of
``Zsigmondy'', ``primitive divisor'', ``support problem'', ``linear recurrence'', ``Ailon'',
``synchronization cost'', ``Hadamard matrix'', or ``preconditioning''.

More importantly, the mathematical objects differ as follows.

\begin{longtable}{@{}p{0.30\textwidth}p{0.31\textwidth}p{0.32\textwidth}@{}}
\toprule
New object in this report & Nearest object in the unified report & Essential difference \\
\midrule
\endfirsthead
\toprule
New object in this report & Nearest object in the unified report & Essential difference \\
\midrule
\endhead
Same-direction quotient
$Q_{\mathbf a}(n)$ & Rational secants between distinct lattice nodes & Here one fixes a
single multiplicative direction and raises both rational ratios to the same layer $n$;
cyclotomic divisibility, not interpolation geometry, is central. \\
Denominator saturation & Denominators of transferred convergents/secants & The new rate is
an exact asymptotic $n\log A+o(n)$ obtained from a homogeneous generalized gcd theorem in
every direction. \\
Recurrence signature of $g_{2,M}$ & General references to BCZ-type gcd bounds & The 2026
classification detects integrality by $C$-finiteness and rules out every nonperiodic common
LDS under independence. \\
Support defects & Prime valuations in determinant entries & The support theorem ignores
multiplicity and converts almost-all-prime order inclusion directly into a power relation. \\
Capture cost $\kappa(n)$ & $p$-adic transport ceilings & Capture cost is a multiplicative
order quotient: the exact least exponent multiplier needed to reproduce a whole cleaned
cyclotomic gap. \\
Three-gap triangle & Adjacent smooth-node secants & The new statement is an exact pairwise
gcd identity and an exponential lcm computation for the three input gaps. \\
Hadamard sign preconditioning & Entrywise/Hadamard powers of matrices & A Hadamard matrix
acts by orthogonal row mixing; the report proves analytic exterior-power neutrality and
$p$-adic Smith neutrality away from primes dividing its order. \\
Formalization plan & Informal proof-roadmap sections & The Lean plan splits elementary,
computationally certifiable arithmetic from three explicit deep theorem interfaces. \\
\bottomrule
\end{longtable}

Only the minimal inherited facts listed in Section~\ref{sec:setup} were reused.  They are
necessary to connect the new invariants to the original conjecture and were not reproved at
the scale of the source report.

\section{Dependency graph and trust boundary}
\label{app:dependencies}

The logical flow is:
\[
\begin{array}{c}
\text{Corvaja--Zannier generalized gcd theorem}
\\[0.3em]\Downarrow\\[-0.2em]
\text{homogeneous gcd estimate}
\\[0.3em]\Downarrow\\[-0.2em]
\text{directional denominator saturation}
\\[0.3em]\Downarrow\\[-0.2em]
\begin{matrix}
\text{fixed-support criteria}&\quad&
\text{capture divergence}&\quad&
\text{positive-mass divergence}
\end{matrix}
\end{array}
\]
Separately,
\[
\text{Corrales--Schoof support theorem}
\Longrightarrow
\text{support-inclusion certificate},
\]
and
\[
\text{Nguyen-Dang recurrence theorem}
\Longrightarrow
\text{recurrence/LDS certificate}.
\]
The gap triangle, exact capture formula, ray covariance, core collapse, numerical
certificates, and Hadamard neutrality are elementary and do not depend on any unproved
conjecture.  The only unresolved implication is from the original simultaneous integrality
hypotheses to one of the displayed certificates.

\section{Proof details omitted from the main flow}
\label{app:proof-details}

\subsection{Why multiplicative independence survives a fixed multiplier}

If $s=r^x$ with $r>1$ and $x$ is a hypothetical noninteger solution, then $r$ and $s^k$ are
multiplicatively independent for every fixed nonzero integer $k$.  Indeed, a relation
$r^u=(s^k)^v$ gives $u=xkv$, hence $x=u/(kv)\in\Q$; the rational-exponent lemma then makes
$x$ integral.  This observation is used every time the homogeneous gcd theorem is applied to
$C^{kn}-D^{kn}$.

\subsection{The local lcm formula}

Let $t_i$ be the order of $s$ modulo the prime-power factors of $\widetilde U_n$.  The order
modulo the full cleaned modulus is $t=\operatorname{lcm}_i t_i$.  For every prime $q$,
\[
 v_q\!\left(\frac{t}{\gcd(t,n)}\right)
 =\max\{0,\max_i v_q(t_i)-v_q(n)\}
 =\max_i v_q\!\left(\frac{t_i}{\gcd(t_i,n)}\right).
\]
This proves
\[
 \frac{t}{\gcd(t,n)}
 =\operatorname{lcm}_i\frac{t_i}{\gcd(t_i,n)}.
\]

\subsection{A corrected radical criterion}

A quantitatively safe version of Theorem~\ref{thm:criteria-menu}(e) is the following.  If
along an unbounded sequence
\[
 \log\rad(\Delta_n)=o\!\left(\frac{n}{\log n}\right)
\]
and every $v_p(\Delta_n)=O(\log n)$ with an absolute implied constant, then
$\log\Delta_n=o(n)$ and $x$ is integral.  The factor $1/\log n$ cannot be omitted without a
stronger multiplicity estimate.

\section{Research provenance and literature notes}
\label{app:provenance}

The attached unified report was read as the baseline, not merely used as a source of the
problem statement.  The recent-literature search then focused on gcd sequences, support
problems, primitive divisors, and verifiable public records of the breakthroughs mentioned in
the prompt.

Nguyen-Dang's June 2026 preprint is the only post-2025 theorem used substantively.  Its main
recurrence and common-LDS statements are imported exactly as external results; its recent
status is explicitly marked.  The denominator theorem relies on the established
Corvaja--Zannier generalized gcd estimate.  The support criterion relies on the 1997 theorem
of Corrales-Rodrig\'a\~nez and Schoof.

For the news audit, an August 2026 arXiv paper gives a self-contained account of the announced
three-dimensional Jacobian counterexample.  By contrast, public Hadamard construction
records available during this work still listed order $668$ as unresolved.  This discrepancy
is why the report uses Hadamard matrices only conditionally and does not rely on the prompt's
claim that all twelve missing orders below $2000$ have been constructed.

\begin{thebibliography}{99}

\bibitem{BCZ2003}
Y.~Bugeaud, P.~Corvaja, and U.~Zannier,
\newblock An upper bound for the G.C.D. of $a^n-1$ and $b^n-1$,
\newblock \emph{Mathematische Zeitschrift} \textbf{243} (2003), 79--84,
\newblock \url{https://doi.org/10.1007/s00209-002-0449-z}.

\bibitem{CatiPasechnik2024}
P.~Cati and D.~V. Pasechnik,
\newblock A database of constructions of Hadamard matrices,
\newblock arXiv:2411.18897, 2024; revised construction tables consulted in 2026,
\newblock \url{https://arxiv.org/abs/2411.18897}.

\bibitem{CorralesSchoof1997}
C.~Corrales-Rodrig\'a\~nez and R.~Schoof,
\newblock The support problem and its elliptic analogue,
\newblock \emph{Journal of Number Theory} \textbf{64} (1997), 276--290,
\newblock \url{https://doi.org/10.1006/jnth.1997.2114}.

\bibitem{CorvajaZannier2005}
P.~Corvaja and U.~Zannier,
\newblock A lower bound for the height of a rational function at $S$-unit points,
\newblock \emph{Monatshefte f\"ur Mathematik} \textbf{144} (2005), 203--224,
\newblock arXiv:math/0311030,
\newblock \url{https://arxiv.org/abs/math/0311030}.

\bibitem{Gao2026}
S.~Gao,
\newblock Counterexamples to the Jacobian conjecture in dimensions greater than two,
\newblock arXiv:2608.00222, 2026,
\newblock \url{https://arxiv.org/abs/2608.00222}.

\bibitem{NguyenDang2026}
K.-H.~Nguyen-Dang,
\newblock On the sequence $\gcd(a^n-1,b^n-1)$,
\newblock arXiv:2606.07959v1, 6 June 2026,
\newblock \url{https://arxiv.org/abs/2606.07959}.

\bibitem{SageHadamard}
The SageMath Developers,
\newblock \texttt{sage.combinat.matrices.hadamard\_matrix}, construction database and
existence queries,
\newblock source consulted 22 August 2026,
\newblock \url{https://github.com/sagemath/sage/blob/develop/src/sage/combinat/matrices/hadamard_matrix.py}.

\bibitem{Zsigmondy1892}
K.~Zsigmondy,
\newblock Zur Theorie der Potenzreste,
\newblock \emph{Monatshefte f\"ur Mathematik und Physik} \textbf{3} (1892), 265--284.

\end{thebibliography}

\end{document}
